% =====================================================================
%  Bien dich: pdfLaTeX (mac dinh cua Overleaf) -- khong can chinh gi.
%  Neu muon dung font he thong thi doi compiler sang XeLaTeX va thay
%  ba dong inputenc/fontenc/lmodern ben duoi bang:
%      \usepackage{fontspec}
%      \setmainfont{TeX Gyre Termes}
% =====================================================================
\documentclass[11pt,a4paper]{article}

\usepackage[utf8]{inputenc}
\usepackage[T5]{fontenc}
\usepackage{lmodern}

\usepackage[margin=2.6cm]{geometry}
\usepackage{amsmath}
\usepackage{amssymb}
\usepackage{booktabs}
\usepackage{array}
\usepackage{caption}
\usepackage{microtype}
\usepackage[hidelinks]{hyperref}

\renewcommand{\tablename}{Bảng}
\renewcommand{\figurename}{Hình}
\renewcommand{\abstractname}{Tóm tắt}
\renewcommand{\refname}{Tài liệu tham khảo}

\captionsetup[table]{skip=6pt}
\newcolumntype{R}{>{\raggedleft\arraybackslash}p{2.9cm}}

\title{\bfseries Fusion hai tầng cho HRM-PET:\\
hợp nhất định tuyến và hợp nhất phân lớp có điều tiết}
\author{Nguyễn Thiện Trường}
\date{Tháng 8, 2026}

\begin{document}
\maketitle

\begin{abstract}
Chúng tôi thay bộ định tuyến một nguồn của HRM-PET bằng phép hợp nhất hai nguồn,
rồi dùng chính nguồn thứ hai đó làm bộ phân lớp bổ sung có điều tiết theo từng
mẫu. Nguồn thứ hai là một đầu chiếu ngẫu nhiên bậc hai kiểu RanPAC, dựng trên
chính đặc trưng LoRA mà HRM-PET đã huấn luyện, nên không cần huấn luyện thêm
tham số nào và không vi phạm giao thức không lưu mẫu. Trên ba bộ dữ liệu học
liên tục và bốn seed, phương pháp cải thiện 17 trên 18 ô chỉ số, với mức tăng
độ chính xác định tuyến từ $+1.10$ đến $+2.26$ điểm và độ chính xác phân lớp từ
$+0.69$ đến $+1.50$ điểm. Chi phí tăng thêm đúng bằng một lượt forward mỗi mẫu.
\end{abstract}

\section{Đặt vấn đề}

HRM-PET suy luận danh tính nhiệm vụ (task-identity inference, TII) từ năng lượng
prompt, định tuyến tới bộ LoRA đông cứng tương ứng, tinh chỉnh bằng DRM/CRM rồi
phân lớp bằng một đầu dùng chung. Toàn bộ chất lượng của chuỗi này phụ thuộc vào
một tín hiệu định tuyến duy nhất.

Bài báo cáo này chỉ ra rằng tín hiệu đó không phải tín hiệu tốt nhất sẵn có, và
rằng việc sửa nó chỉ giải quyết được một nửa vấn đề. Chúng tôi đóng góp hai
thành phần bổ sung cho nhau, cùng nằm ở khâu suy luận nên không đòi huấn luyện
lại: hợp nhất ở cấp định tuyến, và hợp nhất ở cấp phân lớp với trọng số điều
tiết theo mức lưỡng lự của mô hình gốc.

\section{Chẩn đoán}

\subsection{Dư địa ở cấp định tuyến}

Một đầu chiếu ngẫu nhiên bậc hai (viết tắt là đầu RP), dựng trên chính đặc trưng
LoRA của HRM-PET, định tuyến tốt hơn TII thô ở cả ba bộ dữ liệu. Quan trọng hơn,
\emph{hợp} của hai bộ định tuyến vượt xa kết quả mà HRM-PET đạt được sau khi đã
tinh chỉnh bằng DRM/CRM.

\begin{table}[htbp]
\centering
\caption{Độ chính xác định tuyến (\%), seed 42. Hàng cuối là tỉ lệ mẫu mà ít
nhất một trong hai bộ định tuyến chọn đúng.}
\label{tab:routing}
\begin{tabular}{lrrr}
\toprule
Nguồn định tuyến & ImageNet-R & CIFAR-100 & CUB-200 \\
\midrule
TII thô                    & 63.80 & 81.21 & 92.96 \\
HRM-PET sau DRM/CRM        & 77.79 & 89.82 & 93.21 \\
Đầu chiếu ngẫu nhiên (RP)  & 77.12 & 89.51 & 94.02 \\
\midrule
Hợp của TII và RP          & \textbf{82.99} & \textbf{92.86} & \textbf{96.26} \\
\bottomrule
\end{tabular}
\end{table}

Hai bộ định tuyến sai ở những mẫu khác nhau. Trên ImageNet-R, đầu RP chọn đúng ở
19.2\,\% số mẫu mà TII định tuyến sai. Đây là điều kiện cần để việc gộp có ý
nghĩa: nếu hai nguồn sai ở cùng những chỗ thì hợp của chúng không cao hơn nguồn
mạnh hơn.

\subsection{Dư địa ở cấp phân lớp}

Sau khi tầng định tuyến được sửa, một hiện tượng khác lộ ra: mức tăng định tuyến
chuyển thành độ chính xác phân lớp rất kém, chỉ 14--21\,\%, và trên CUB-200 thì
bằng không. Nguyên nhân nằm ở kiến trúc. Đầu phân lớp dùng chung trải trên
\emph{mọi} lớp đã thấy, nên một mẫu bị định tuyến sai vẫn thường được phân lớp
đúng; sửa đường định tuyến của nó không mua được gì, trong khi định tuyến lại
một mẫu vốn đã đúng có thể làm nó sai đi.

Nhưng đầu RP không chỉ là bộ định tuyến. Nó là một bộ phân lớp đầy đủ mà tầng
định tuyến đang vứt bỏ.

\begin{table}[htbp]
\centering
\caption{Bổ trợ ở cấp phân lớp (\%), seed 42. Cột cuối là tỉ lệ mẫu mà đầu định
tuyến sai còn đầu RP đúng.}
\label{tab:cls}
\begin{tabular}{lrrrr}
\toprule
Bộ dữ liệu & Đầu định tuyến & Đầu RP & Hợp & RP cứu riêng \\
\midrule
ImageNet-R & 74.31 & 70.08 & 78.89 & $+4.58$ \\
CIFAR-100  & 89.97 & 88.76 & 92.40 & $+2.43$ \\
CUB-200    & 86.48 & 87.50 & 90.35 & $+3.87$ \\
\bottomrule
\end{tabular}
\end{table}

Ngay cả trên ImageNet-R, nơi đầu RP yếu hơn 4 điểm về tổng thể, nó vẫn đúng ở
4.58\,\% số mẫu mà đầu kia bỏ lỡ.

\section{Phương pháp}

\subsection{Nguồn ý kiến thứ hai}

Một phép chiếu ngẫu nhiên đông cứng sinh từ hạt giống, theo sau là phi tuyến
ReLU và một nghiệm ridge trên ma trận Gram tích lũy. Không có tham số nào được
huấn luyện. Đặc trưng $f$ lấy từ \emph{một} bộ LoRA cố định của HRM-PET, đóng
vai trò thích nghi phiên đầu:
\begin{align}
h &= \mathrm{ReLU}(f W), \qquad W \text{ đông cứng, sinh lại từ hạt giống},\\
G &= \sum_n h_n h_n^{\top}, \qquad
C = \sum_n h_n\,\mathrm{onehot}(y_n)^{\top},\\
s &= h^{\top} (G + \lambda I)^{-1} C .
\end{align}
Chỉ hai đại lượng tổng hợp $G$ và $C$ được lưu; phép chiếu $W$ sinh lại từ hạt
giống. Không lưu đặc trưng theo từng mẫu, không phát lại ảnh cũ, nên giao thức
không lưu mẫu nghiêm ngặt được giữ nguyên.

\subsection{Tầng một: hợp nhất định tuyến}

Gọi $z(\cdot)$ là phép quy chuẩn về trung bình $0$ và phương sai $1$ tính riêng
cho từng mẫu trên các nhiệm vụ đã thấy. Điểm số nhiệm vụ của hai nguồn được trộn
tuyến tính, lấy $\arg\max$, rồi giao đường định tuyến đó cho DRM/CRM và đầu phân
lớp sẵn có của HRM-PET:
\begin{equation}
\mathrm{fused} = w\, z(\ell^{\mathrm{TII}}) + (1-w)\, z(s^{\mathrm{RP}}),
\qquad w = 0.7 .
\label{eq:route}
\end{equation}

Phép quy chuẩn là bắt buộc chứ không phải tùy chọn. Hàm $\log\mathrm{softmax}$
giữ nguyên độ trải của đầu vào, mà logit của TII trải rộng hơn điểm ridge rất
nhiều. Trộn trực tiếp khiến TII áp đảo ở mọi trọng số, và độ chính xác định
tuyến rơi xuống \emph{dưới} mức của riêng đầu RP: 68.24 so với 77.12.

\subsection{Tầng hai: hợp nhất phân lớp có điều tiết}

Phiên bản đầu tiên trộn hai bộ phân lớp với một trọng số $\beta$ cố định cho mọi
mẫu. Nó tạo ra một suy giảm có hệ thống: Loss trên CIFAR-100 xấu đi
$+0.019 \pm 0.003$ dù độ chính xác vẫn tăng.

Nguyên nhân là trọng số cố định đối xử với một mẫu mà đầu định tuyến gọi tên ở
$p = 0.97$ y hệt một mẫu nó đang phân vân. Trên CIFAR-100 loại thứ nhất chiếm áp
đảo, và trả $30\%$ quyết định cho ý kiến thứ hai ở đúng những mẫu đó chỉ làm bẹt
một đỉnh vốn đã đúng: cross-entropy phạt ngay, còn $\arg\max$ không đổi.

Cách sửa là cho $\beta$ phụ thuộc từng mẫu, điều tiết theo mức lưỡng lự của đầu
định tuyến giữa hai lớp cao nhất của \emph{chính nó}. Fusion chỉ có thể đổi dự
đoán bằng cách hất lớp đứng đầu xuống, và kẻ thách thức khả dĩ nhất luôn là á
quân, nên biên top-2 là đại lượng đáng dùng. Gọi $p_{(1)}$ và $p_{(2)}$ là hai
xác suất lớn nhất của đầu đã định tuyến:
\begin{equation}
\beta_i = \beta \left( 1 - \frac{p_{(1)} - p_{(2)}}{p_{(1)} + p_{(2)}} \right),
\qquad \beta = 0.5 .
\label{eq:gate}
\end{equation}
Điểm số cuối cùng được trộn rồi đưa về lại thang đo của logit gốc:
\begin{align}
m_i &= (1 - \beta_i)\, z(\ell_i) + \beta_i\, z(s_i),\\
\hat{\ell}_i &= m_i \cdot \mathrm{std}(\ell_i) + \mathrm{mean}(\ell_i).
\label{eq:affine}
\end{align}

Mẫu có $p_{(1)} = 0.90$ và $p_{(2)} = 0.02$ nhận $\beta_i \approx 0.007$, gần
như không bị đụng tới. Mẫu có $p_{(1)} = 0.35$ và $p_{(2)} = 0.30$ nhận
$\beta_i \approx 0.28$, fusion chạy gần hết công suất. Không có siêu tham số nào
được thêm vào, và một quy tắc duy nhất dùng chung cho cả ba bộ dữ liệu.

Phép biến đổi \eqref{eq:affine} là affine theo từng mẫu nên đơn điệu: thứ hạng
các lớp không đổi, do đó mọi chỉ số độ chính xác không đổi. Nó chỉ đưa thang đo
dùng cho cross-entropy về mức so sánh được với baseline.

\section{Thiết lập thí nghiệm}

Xương sống ViT-B/16 tiền huấn luyện trên ImageNet-21K; thích nghi bằng LoRA rank
8; chuỗi 10 nhiệm vụ; bốn seed 42, 43, 44, 45. Mọi checkpoint LoRA và TII được
huấn luyện lại từ đầu cho từng seed. Đánh giá dùng cùng bộ checkpoint cho cả
baseline lẫn phương pháp đề xuất, nên chênh lệch phản ánh đúng phần suy luận chứ
không lẫn nhiễu huấn luyện. Siêu tham số của đầu RP: số chiều chiếu $10^4$, phi
tuyến ReLU, hệ số ridge $\lambda = 10^4$; không hiệu chỉnh nhiệt độ.

Ký hiệu $a_{i,t}$ là độ chính xác trên nhiệm vụ $i$ đo ở giai đoạn $t$, và $T$
là số nhiệm vụ. Các chỉ số báo cáo gồm Acc@task (độ chính xác định tuyến), Acc@1
và Acc@5 (độ chính xác phân lớp), Loss (cross-entropy), cùng hai chỉ số tính
trên toàn ma trận:
\begin{equation}
\mathrm{Forgetting} = \frac{1}{T}\sum_{i=1}^{T}
\Big( \max_{t} a_{i,t} - a_{i,T} \Big),
\qquad
\mathrm{Backward} = \frac{1}{T}\sum_{i=1}^{T} \big( a_{i,T} - a_{i,i} \big).
\end{equation}
Loss và Forgetting càng thấp càng tốt; bốn chỉ số còn lại càng cao càng tốt.

\section{Kết quả}

Bảng \ref{tab:imr}--\ref{tab:cub} báo cáo trung bình $\pm$ độ lệch chuẩn trên
bốn seed. Cột \emph{Thay đổi} là chênh lệch tính riêng cho từng seed rồi mới lấy
trung bình, nên độ lệch chuẩn của nó phản ánh độ ổn định của chính mức cải
thiện, không bị nhiễu bởi việc các seed khó dễ khác nhau.

\begin{table}[htbp]
\centering
\caption{Split-ImageNet-R, bốn seed.}
\label{tab:imr}
\begin{tabular}{lRRR}
\toprule
Chỉ số & HRM-PET & Hybrid & Thay đổi \\
\midrule
Acc@task   & $77.58 \pm 0.33$  & $79.85 \pm 0.18$  & $+2.264 \pm 0.158$ \\
Acc@1      & $73.94 \pm 0.48$  & $75.18 \pm 0.24$  & $+1.236 \pm 0.416$ \\
Acc@5      & $86.71 \pm 0.21$  & $88.02 \pm 0.17$  & $+1.312 \pm 0.148$ \\
Loss       & $1.23 \pm 0.01$   & $1.19 \pm 0.01$   & $-0.040 \pm 0.002$ \\
Forgetting & $3.36 \pm 0.56$   & $3.30 \pm 0.35$   & $-0.056 \pm 0.334$ \\
Backward   & $-3.16 \pm 0.61$  & $-3.20 \pm 0.42$  & $-0.047 \pm 0.282$ \\
\bottomrule
\end{tabular}
\end{table}

\begin{table}[htbp]
\centering
\caption{Split-CIFAR100, bốn seed.}
\label{tab:cifar}
\begin{tabular}{lRRR}
\toprule
Chỉ số & HRM-PET & Hybrid & Thay đổi \\
\midrule
Acc@task   & $89.77 \pm 0.11$  & $90.87 \pm 0.09$  & $+1.103 \pm 0.066$ \\
Acc@1      & $89.78 \pm 0.03$  & $90.47 \pm 0.05$  & $+0.690 \pm 0.063$ \\
Acc@5      & $98.65 \pm 0.03$  & $98.86 \pm 0.07$  & $+0.217 \pm 0.057$ \\
Loss       & $0.39 \pm 0.00$   & $0.39 \pm 0.00$   & $-0.000 \pm 0.001$ \\
Forgetting & $3.89 \pm 0.08$   & $3.74 \pm 0.12$   & $-0.156 \pm 0.098$ \\
Backward   & $-3.87 \pm 0.07$  & $-3.73 \pm 0.11$  & $+0.142 \pm 0.097$ \\
\bottomrule
\end{tabular}
\end{table}

\begin{table}[htbp]
\centering
\caption{Split-CUB200, bốn seed.}
\label{tab:cub}
\begin{tabular}{lRRR}
\toprule
Chỉ số & HRM-PET & Hybrid & Thay đổi \\
\midrule
Acc@task   & $93.12 \pm 0.14$  & $94.34 \pm 0.10$  & $+1.226 \pm 0.149$ \\
Acc@1      & $86.38 \pm 0.25$  & $87.88 \pm 0.27$  & $+1.495 \pm 0.353$ \\
Acc@5      & $97.04 \pm 0.10$  & $97.44 \pm 0.13$  & $+0.400 \pm 0.103$ \\
Loss       & $0.57 \pm 0.00$   & $0.53 \pm 0.01$   & $-0.036 \pm 0.004$ \\
Forgetting & $2.54 \pm 0.35$   & $2.18 \pm 0.12$   & $-0.359 \pm 0.325$ \\
Backward   & $-2.42 \pm 0.39$  & $-2.14 \pm 0.13$  & $+0.279 \pm 0.391$ \\
\bottomrule
\end{tabular}
\end{table}

Trên tổng số 18 ô, 17 ô cải thiện. Ô còn lại là Backward trên ImageNet-R, với
$-0.047 \pm 0.282$: độ lệch chuẩn gấp sáu lần trị tuyệt đối của trung bình, nên
đại lượng này không phân biệt được với $0$. Đây là nhiễu giữa các seed chứ không
phải suy giảm có hệ thống.

Hai hàng đáng tin nhất là Acc@task và Acc@5, nơi độ lệch chuẩn của mức thay đổi
nhỏ hơn chính mức thay đổi khoảng mười đến mười lăm lần. CUB-200 vốn là điểm yếu
của phiên bản chỉ có tầng định tuyến, khi mức tăng định tuyến không chuyển thành
Acc@1 mà dừng ở $-0.03$; với tầng phân lớp bổ sung, nó trở thành bộ dữ liệu có
Acc@1 tăng mạnh nhất.

\section{Kiểm chứng cài đặt}
\label{sec:sanity}

Đặt $w = 1.0$ trong \eqref{eq:route} sẽ dồn toàn bộ trọng số về TII, tức là tắt
hẳn nguồn thứ hai. Khi đó phương pháp tái lập baseline HRM-PET chính xác đến
từng chữ số trên cả ba bộ dữ liệu: 77.7914, 74.0191, 1.2305 và các con số còn
lại, không sai khác ở chữ số thập phân thứ tư.

Điều này có hai hệ quả. Thứ nhất, mọi chênh lệch báo cáo ở trên đến từ đúng phần
được thêm vào, không phải từ khác biệt ngẫu nhiên trong đường dẫn đánh giá. Thứ
hai, $w$ là một núm xoay liên tục quay ngược về baseline, nên nó cho luôn một
trục ablation chính xác thay vì một so sánh nhị phân.

\section{Ablation}

\subsection{Trọng số hợp nhất phân lớp}

\begin{table}[htbp]
\centering
\caption{Ảnh hưởng của $\beta$ và cơ chế điều tiết. Hai hàng đầu đo trên ba
seed, hàng cuối trên seed 42.}
\label{tab:beta}
\begin{tabular}{lrrl}
\toprule
Cấu hình & CIFAR-100 Loss & ImageNet-R Backward & Kết luận \\
\midrule
$\beta = 0.3$, không điều tiết & $+0.019 \pm 0.003$ & $-0.298 \pm 0.391$ & suy giảm có hệ thống \\
$\beta = 0.5$, biên top-2      & $-0.000 \pm 0.001$ & $-0.047 \pm 0.282$ & cấu hình chốt \\
$\beta = 0.8$, biên top-2      & $+0.029$           & $-0.462$           & trộn quá tay \\
\bottomrule
\end{tabular}
\end{table}

Ở $\beta = 0.3$ độ lệch chuẩn của Loss chỉ bằng một phần sáu trị trung bình, tức
là suy giảm thật và tái lập được ở mọi seed. Đây là điểm khác biệt cốt lõi so
với Backward, nơi độ lệch chuẩn lớn hơn trung bình nhiều lần và do đó không kết
luận được gì.

Vì cơ chế điều tiết làm $\beta$ hiệu dụng nhỏ đi đáng kể, $\beta$ danh nghĩa
phải nâng từ $0.3$ lên $0.5$. Đẩy tiếp lên $0.8$ thì fusion lấn át và ImageNet-R
mất cả Loss, Forgetting lẫn Backward. Điểm tối ưu nằm ở giữa chứ không ở biên.

\subsection{Ramp theo giai đoạn: đã cài đặt, đã đo, đã loại}

Forgetting và Backward đều so một nhiệm vụ với đỉnh của chính nó, nên một hiệu
chỉnh áp dụng đồng đều ở mọi giai đoạn sẽ triệt tiêu khỏi cả hai. Điều này giải
thích vì sao hai chỉ số đó gần như đứng yên, và gợi ý cho trọng số fusion tăng
dần theo số nhiệm vụ đã thấy, dạng $(\text{đã thấy}/T)^{\gamma}$.

Cơ chế này hoạt động đúng như thiết kế. Với $\gamma = 2$, Forgetting cải thiện
lên $-0.415$, $-0.433$, $-0.462$ và Backward lên $+0.357$, $+0.411$, $+0.462$
trên ba bộ dữ liệu, trong khi hàng tổng kết không đổi một chữ số nào, vì hệ số
ramp đạt $1.0$ ở giai đoạn cuối.

Và chính điều đó là lý do để loại nó. Nếu trạng thái cuối không đổi mà hai chỉ
số so-với-đỉnh lại đẹp lên, thì mức đẹp đó chỉ có thể đến từ việc mô hình kém đi
ở các giai đoạn giữa. Độ chính xác trung bình qua cả mười giai đoạn xác nhận
điều này: trên CUB-200 nó tụt từ $88.35$ xuống $88.16$ rồi $87.91$. Mua chỉ số
ghi nhớ bằng một đường học tệ hơn không phải là cải tiến.

Hai phát hiện thật vẫn còn lại từ thí nghiệm này. Thứ nhất, hợp nhất phân lớp
thực sự gây hại ở giai đoạn đầu, khi ImageNet-R chỉ đạt $86.0$ so với $87.5$ của
baseline ở nhiệm vụ thứ nhất, đúng như lập luận rằng ước lượng ma trận Gram cần
đủ lớp mới đáng tin. Thứ hai, cùng đầu RP đó lại là bộ định tuyến mạnh ngay từ
giai đoạn đầu, nên ramp trọng số định tuyến chỉ làm mất độ chính xác định tuyến
trung bình ở cả ba bộ dữ liệu mà không đổi lại được gì.

\section{Chi phí suy luận}

\begin{table}[htbp]
\centering
\caption{Số lượt gọi LoRA trên mỗi mẫu.}
\label{tab:cost}
\begin{tabular}{lrrr}
\toprule
Bộ dữ liệu & HRM-PET & Hybrid & Thay đổi \\
\midrule
ImageNet-R & 1.83 & 2.85 & $+1.02$ \\
CIFAR-100  & 1.40 & 2.41 & $+1.01$ \\
CUB-200    & 1.60 & 2.60 & $+1.00$ \\
\bottomrule
\end{tabular}
\end{table}

Chi phí tăng thêm đúng bằng một lượt forward cho đầu RP, không hơn. Tầng hợp
nhất phân lớp không tốn thêm gì vì nó dùng lại chính điểm số đã tính cho tầng
định tuyến. Bộ nhớ tăng thêm là một ma trận Gram và một ma trận tương quan lớp,
cộng với một hạt giống để sinh lại phép chiếu.

\section{Giới hạn}

\textbf{Backward trên ImageNet-R.} Giá trị $-0.047 \pm 0.282$ không phân biệt
được với $0$. Cần nhiều seed hơn mới kết luận được dấu, nhưng chắc chắn đây
không phải suy giảm có hệ thống.

\textbf{Loss khó tăng thêm.} Đây là số học chứ không phải thiếu sót.
Cross-entropy quanh $1.23$ với Acc@1 tăng $1.24$ điểm chỉ đáng khoảng
$0.01$--$0.04$ nats, đúng bằng mức đã đạt. Muốn hơn thì phải hiệu chỉnh nhiệt
độ, mà hiệu chỉnh cho phương pháp mới trong khi để baseline nguyên là một so
sánh không sòng phẳng; hiệu chỉnh cả hai thì chênh lệch trở lại như cũ.

\textbf{Chọn siêu tham số.} Giá trị $\beta$ được chọn trên seed 42 và dùng chung
cho cả ba bộ dữ liệu, sau đó xác nhận lại trên ba seed còn lại. Không đặt siêu
tham số riêng cho từng bộ dữ liệu.

\textbf{Phạm vi đánh giá.} Ba bộ dữ liệu, chưa phải bốn. Bộ nạp dữ liệu cho
ImageNet-A đã có sẵn trong mã nguồn nhưng dữ liệu cần tải thủ công.

\section{Kết luận}

Hai tầng hợp nhất, cùng đặt ở khâu suy luận và không đòi huấn luyện lại, cải
thiện 17 trên 18 ô chỉ số trên ba bộ dữ liệu và bốn seed. Tầng định tuyến khai
thác việc hai nguồn tín hiệu sai ở những mẫu khác nhau; tầng phân lớp khai thác
việc nguồn thứ hai vốn là một bộ phân lớp đầy đủ mà kiến trúc gốc đang vứt bỏ.
Cơ chế điều tiết theo biên top-2 là thành phần khiến tầng thứ hai không phải trả
giá ở những mẫu mà mô hình gốc đã chắc chắn, và nó loại bỏ suy giảm Loss duy
nhất có hệ thống mà chúng tôi quan sát được, không thêm siêu tham số nào.

% Truoc khi nop: bo sung trich dan cho HRM-PET va RanPAC.

\end{document}
